\clearpage\SourceLeaf{253}
\par\bigskip\Needspace{9\baselineskip}
\pdfbookmark[0]{Further Development of the Investigations of Pure Equations}{aequationes-purae}
\begin{center}\large\bfseries FURTHER DEVELOPMENT OF THE INVESTIGATIONS\par\medskip OF PURE EQUATIONS.\end{center}
\BookRule{.12\linewidth}
\renewcommand{\BookRunningTitle}{Further Investigations of Pure Equations}
\renewcommand{\Article}[1]{\par\medskip\Needspace{4\baselineskip}\pdfbookmark[1]{Art. #1}{purae-art-#1}\centerline{\textbf{#1.}}\nobreak\smallskip}

\Article{1}
The method by which, in Article~360 of the \emph{Disquiss.\ Arithm.}, we showed how to solve the equation \mbox{$x^n-1=0$} forms a most fruitful and important theory, of which we could touch on only the first elements in that work. We therefore hope that mathematicians will welcome our taking up this subject here again, treating more fully what was only briefly sketched, in part with the proofs omitted, and pursuing in greater depth the developments made since that time.

The exponent $n$ is assumed to be a prime number, and the number $n-1$ to be resolved into the factors $\alpha\times\beta\times\gamma$; further, we denote by $g$ any primitive root modulo $n$. Let $r$ denote an arbitrary root of the equation \mbox{$x^n-1=0$}, and $R$ an arbitrary root of the equation \mbox{$x^\beta-1=0$}. Thus, denoting the circumference of the circle whose radius $=1$ by $P$, and the imaginary quantity $\mathord{\surd}{-1}$ by $i$, all the roots of the equation \mbox{$x^\beta-1=0$}, or all the values of $R$, will be represented by the formula
\[
\cos\frac{kP}{\theta}+i\sin\frac{kP}{\theta}
\]
where $k$ ranges over the integers $0,1,2,3\ldots\beta-1$. Further, all powers of any root $R$ are evidently themselves roots, and if $R$ is the root corresponding to a value of $k$ coprime to $\beta$, all the powers $R^0,R,R^2,R^3\ldots R^{\beta-1}$ are distinct from one another and thus exhaust the entire collection of roots. In this case we shall call $R$ a \emph{proper} root of the equation \mbox{$x^\beta-1=0$}; conversely, a root $R$
\SourceLeaf{254}
corresponding to a value of $k$ that is not coprime to $\beta$ will be called \emph{improper}. It is readily seen that, if $\delta$ is the greatest common divisor of the numbers $k$ and $\beta$, then $R^{\frac\beta\delta .}=1$, while all the powers $R^0,R,R^2,R^3\ldots R^{\frac\beta\delta-1}$ are distinct from one another, so that $R$ is a proper root of the equation $x^{\frac\beta\delta}-1=0$. The same holds for the equation \mbox{$x^n-1=0$}, but all its roots are necessarily proper except for the root $1$.

\Needspace{20\baselineskip}\Article{2}
With these preliminaries, our investigation will chiefly concern functions of the following form, consisting of $\beta\gamma$ terms,
\[
r+Rr^{g^\alpha}+R^2r^{g^{2\alpha}}+R^3r^{g^{3\alpha}}\ldots+R^{\beta\gamma-1}r^{g^{\alpha\beta\gamma-\alpha}}
\]
which for brevity we shall denote by the symbol $[r,R]$. The individual terms of such an expression are products of powers of $r$ and powers of $R$; the exponents of the former form a geometric progression, and those of the latter an arithmetic progression. The exponents
\[
1,\quad g^\alpha,\quad g^{2\alpha},\quad g^{3\alpha}\ldots\ldots g^{\alpha\beta\gamma-\alpha}
\]
are all incongruent to one another modulo $n$, and hence those powers of $r$ are distinct from one another; continued further, however, they would begin the same series again, since $g^{\alpha\beta\gamma}\equiv1(\mathrm{mod.}\,n)$, and hence $r^{g^{\alpha\beta\gamma}}=r$. The other factors,
\[
1,\quad R,\quad R^2,\quad R^3\dotsfour R^{\beta\gamma-1}
\]
form $\gamma$ equal periods, since $R^\beta=1$, $R^{\beta-1}=R$, etc. It follows that the function $[r,R]$ can also be represented as
\[
\begin{aligned}
&r+r^{g^{\alpha\beta}}+r^{g^{2\alpha\beta}}\ldots\ldots+r^{g^{\alpha\beta\gamma-\alpha\beta}}\\
&{}+R\bigl(r^{g^\alpha}+r^{g^{\alpha\beta+\alpha}}+r^{g^{2\alpha\beta+\alpha}}\ldots\ldots+r^{g^{\alpha\beta\gamma-\alpha\beta+\alpha}}\bigr)\\
&{}+R^2\bigl(r^{g^{2\alpha}}+r^{g^{\alpha\beta+2\alpha}}+r^{g^{2\alpha\beta+2\alpha}}\ldots\ldots+r^{g^{\alpha\beta\gamma-\alpha\beta+2\alpha}}\bigr)\\
&{}+\text{etc.}\\
&{}+R^{\beta-1}\bigl(r^{g^{\alpha\beta-\alpha}}+r^{g^{2\alpha\beta-\alpha}}+r^{g^{3\alpha\beta-\alpha}}\ldots\ldots+r^{g^{\alpha\beta\gamma-\alpha}}\bigr).
\end{aligned}
\]
or, introducing the notation of Article~343 of the Disq.\ Ar.,
\[
[r,R]=(\gamma,1)+R(\gamma,g^\alpha)+R^2(\gamma,g^{2\alpha})\ldots+R^{\beta-1}(\gamma,g^{\alpha\beta-\alpha})
\]

\SourceLeaf{255}
\Needspace{20\baselineskip}\Article{3}
If we take unity for the root $r$, we have
\[
[1,R]=1+R+R^2+R^3\ldots+R^{\beta\gamma-1}=\gamma(1+R+R^2+R^3\ldots+R^{\beta-1})
\]
whose value will be $=\beta\gamma$ if $R$ is also taken to be the root $1$, but $=0$ for every other value of $R$. Conversely, leaving $r$ indeterminate and putting $R=1$, we obtain $[r,1]=r+r^{g^\alpha}+r^{g^{2\alpha}}+r^{g^{3\alpha}}\ldots+r^{g^{\alpha\beta\gamma-\alpha}}$, or, using the notation introduced in the Disq.\ Ar., $[r,1]=(\beta\gamma,1)$; that is, it will consist of a period of $\beta\gamma$ roots, one of which is $r$ itself. Whenever $\alpha=1$, this period will contain all the roots $r,r^2,r^3\ldots r^{n-1}$, with only their order changed.

Note also the following relations, whose reason is evident:
\[
[r,R]=R[r^{g^\alpha},R]=R^2[r^{g^{2\alpha}},R]\quad\text{or generally}\quad=R^k[r^{g^{\alpha k}},R]
\]
where $k$ denotes any positive integer. It follows that the function $[r^m,R]$ either is $=[1,R]$, namely if $m$ is divisible by $n$, or can be reduced in the remaining cases to the form $R^\mu[r^{g^\nu},R]$, with $\nu<\alpha$. For if $m$ is not divisible by $n$, it will be congruent modulo $n$ to some power of $g$, whose exponent is conveniently expressed by $\mathrm{ind.}\,m$, as in the Disq.\ Ar.; thus, putting $\mathrm{ind.}\,m=\lambda\alpha+\nu$, which can evidently be done with $\nu<\alpha$, we obtain $[r^m,R]=[r^{g^{\lambda\alpha+\nu}},R]=R^{-\lambda}[r^{g^\nu},R]$. We must therefore set $\mu=-\lambda$ or, if you want a positive exponent, $\mu\equiv-\lambda(\mathrm{mod.}\,\beta)$.

\Article{4}
\textsc{Theorem.} \emph{If $r'$, like $r$, denotes an arbitrary root of the equation \mbox{$x^n-1=0$}, and $R'$, like $R$, an arbitrary root of the equation \mbox{$x^\beta-1=0$}, the product will be}
\[
\begin{aligned}
&[r,R]\times[r',R']=\\
&\quad[rr',RR']+R[r^{g^\alpha}r',RR']+R^2[r^{g^{2\alpha}}r',RR']\\
&\quad{}+R^3[r^{g^{3\alpha}}r',RR']\ldots+R^{\beta\gamma-1}[r^{g^{\alpha\beta\gamma-\alpha}}r',RR']
\end{aligned}
\]

\emph{Proof.} By multiplying $[r,R]$ by the individual terms of $[r',R']$, the product can be represented in the form
\[
\begin{aligned}
&[r,R]r'+RR'[r^{g^\alpha},R]r'^{g^\alpha}+R^2R'^2[r^{g^{2\alpha}},R]r'^{g^{2\alpha}}\\
&{}+R^3R'^3[r^{g^{3\alpha}},R]r'^{g^{3\alpha}}\ldots\\
&{}+R^{\beta\gamma-1}R'^{\beta\gamma-1}[r^{g^{\alpha\beta\gamma-\alpha}},R]r'^{g^{\alpha\beta\gamma-\alpha}}
\end{aligned}
\]

\SourceLeaf{256}
\Needspace{7\baselineskip}Then collecting the first terms of the individual parts, after duly expanding them, gives $[rr',RR']$; similarly, collecting the second terms gives $R[r^{g^\alpha}r',RR']$, and so on, finally producing the stated form of the product. Q.\ E.\ D.

Moreover, simply interchanging $r,R$ with $r',R'$ shows that the same product can also be put in the form
\[
\begin{aligned}
&[rr',RR']+R'[rr'^{g^\alpha},RR']+R'^2[rr'^{g^{2\alpha}},RR']\\
&{}+R'^3[rr'^{g^{3\alpha}},RR']\dotsfour+R'^{\beta\gamma-1}[rr'^{g^{\alpha\beta\gamma-\alpha}},RR']
\end{aligned}
\]

\Needspace{16\baselineskip}It follows further that, if $r'',r'''$, etc., also denote arbitrary roots of the equation \mbox{$x^n-1=0$}, and $R'',R'''$, etc., arbitrary roots of the equation \mbox{$x^\beta-1=0$}, the product of the functions $[r,R],[r',R'],[r'',R''],[r''',R''']$, etc., however many there are, will equal the sum
\[
\begin{gathered}
\Sigma R'^{k'}R''^{k''}R'''^{k'''}\text{ etc.}\\
{}[rr'^{g^{\alpha k'}}r''^{g^{\alpha k''}}r'''^{g^{\alpha k'''}}\text{ etc.},\ RR'R''R'''\text{ etc.}]
\end{gathered}
\]
For $k',k'',k'''$, etc., we substitute all the numbers $0,1,2,3\ldots\beta\gamma-1$, combined with one another in every possible different way. This gives altogether $\beta^{\mu-1}\gamma^{\mu-1}$ terms, if $\mu$ denotes the number of those functions multiplied together.
